% Claim proof file for row claim_3_1 -- "JNodeProperties".
% Auto-generated stub by scaffold/solving_current_row.py at row start. A
% manager agent dedicated to writing the TeX proof fills in the body below.
% The proof must:
%   - Stay close to the lecture notes' proof if one exists (always search
%     the LN first for a `\begin{proof}` block immediately following the
%     claim; copy / lightly edit when present).
%   - Otherwise use the LN paradigm and cite prior claims by their `ref`.
%   - Be self-contained enough that a mathematician can follow it without
%     re-reading the LN.
%   - Have no `sorry`, `omitted`, or "obvious" shortcuts.
%
% Compiles standalone via the `subfiles` package.

\documentclass[main]{subfiles}

\begin{document}
% ref: claim_3_1 (proof)
% title: JNodeProperties
%
% Cross-references: see claim_<ref>_statement_<title>.tex in this folder
% for the corresponding statement.

% The LN states this as a remark immediately after \ref{not-cdmg}
% (graphs.tex line 84-89) and provides no proof block; the statement is
% the unfolding-only consequence of the structural fields of
% \ref{def-cdmg}. The proof below is constructed from scratch in the
% LN paradigm: every step is one notation-unfolding of \ref{not-cdmg}
% followed by one of the inclusions $E \ins (J \cup V) \x V$,
% $L \ins V \x V$, or the disjointness $J \cap V = \emptyset$ from
% \ref{def-cdmg}.

\begin{proof}
Let $G = (J, V, E, L)$ be a CDMG as in \ref{def-cdmg}. We verify the
three assertions in turn.

\medskip
\noindent\emph{(1) No arrowhead points at any $j \in J$:
$j \hus v \notin G$ for all $v \in G$.}
Suppose for contradiction that $j \in J$ and $j \hus v \in G$ for some
$v$. By \ref{not-cdmg}, $j \hus v \in G$ means that either
$j \hut v \in G$ or $j \huh v \in G$.
\begin{itemize}
  \item If $j \hut v \in G$, then by \ref{not-cdmg} $(v, j) \in E$.
        By \ref{def-cdmg} we have $E \ins (J \cup V) \times V$,
        so the second coordinate of any pair in $E$ lies in $V$; hence
        $j \in V$.
  \item If $j \huh v \in G$, then by \ref{not-cdmg} $(j, v) \in L$.
        By \ref{def-cdmg} we have $L \ins V \times V$, so both
        coordinates of any pair in $L$ lie in $V$; hence $j \in V$.
\end{itemize}
Either way $j \in J \cap V$, contradicting the disjointness of $J$ and
$V$ stipulated in \ref{def-cdmg}.

\medskip
\noindent\emph{(2) Directed edges out of an input node target $V$:
if $j \in J$ and $j \tuh w \in G$, then $w \in V$.}
By \ref{not-cdmg}, $j \tuh w \in G$ means $(j, w) \in E$. Since
$E \ins (J \cup V) \times V$ by \ref{def-cdmg}, the second coordinate
$w$ lies in $V$.

\medskip
\noindent\emph{(3) No two input nodes are adjacent: for all
$j_1, j_2 \in J$, $j_1 \sus j_2 \notin G$.}
Suppose for contradiction that $j_1, j_2 \in J$ and
$j_1 \sus j_2 \in G$. By \ref{not-cdmg}, $j_1 \sus j_2 \in G$ means
one of the three sub-cases $j_1 \tuh j_2 \in G$,
$j_1 \hut j_2 \in G$, or $j_1 \huh j_2 \in G$.
\begin{itemize}
  \item If $j_1 \tuh j_2 \in G$, then $(j_1, j_2) \in E$, which by
        \ref{not-cdmg} is exactly $j_2 \hut j_1 \in G$. In particular
        $j_2 \hus j_1 \in G$, contradicting part (1) applied to
        $j_2 \in J$.
  \item If $j_1 \hut j_2 \in G$, then $j_1 \hus j_2 \in G$ directly
        by the definition of $\hus$ in \ref{not-cdmg}, contradicting
        part (1) applied to $j_1 \in J$.
  \item If $j_1 \huh j_2 \in G$, then $j_1 \hus j_2 \in G$ (again
        directly from \ref{not-cdmg}), contradicting part (1)
        applied to $j_1 \in J$.
\end{itemize}
In each case we reach a contradiction, so $j_1 \sus j_2 \notin G$; by
the adjacency definition immediately following \ref{not-cdmg}, this
is precisely the statement that no two nodes in $J$ are adjacent.
\end{proof}

\end{document}
